%% Citas verificadas contra el PDF. Sin marcadores [VERIFY] pendientes.

\section{CONCLUSIONES Y TRABAJO FUTURO}

Este trabajo ha presentado el diseño de un sistema que estima el riesgo de estafa en el punto
de decisión de compra, a partir de características que el comprador ya tiene delante y sin
acceso privilegiado a la plataforma. La contribución principal es el marco conceptual, no una
implementación completa: la distinción entre riesgo detectado e información insuficiente, la
regla de abstención que la operationaliza y el protocolo de validación que separa la
observabilidad de la eficacia. Tres conclusiones se desprenden del diseño.

La primera es que la arquitectura propuesta permite plantear el desplazamiento de la
evaluación de riesgo desde el juicio del comprador hacia el sistema sin requerir la
cooperación de la plataforma: las ocho características de la Tabla~\ref{tab:senales} se
derivan del contenido que el navegador ya recibió, lo que distingue la propuesta de la línea
dominante en detección de fraude \cite{ref28}. Que ese desplazamiento resulte además útil para
el comprador es lo que el protocolo de validación de la Sección~III-B deberá establecer.

La segunda es que la regla de abstención explícita no es un detalle de implementación sino
una condición de diseño. Un instrumento que presentara como seguro aquello que solo desconoce
trasladaría al comprador un riesgo que no ha evaluado, y penalizaría a comercios legítimos
que simplemente no publicitan sus garantías.

La tercera es que intervenir sobre la información del comprador tiene fundamento teórico y no
solo utilidad práctica: la cantidad de reseñas falsas en equilibrio decrece con la fracción de
consumidores capaces de interpretar correctamente los indicadores de una ficha \cite{ref29}.
De ese resultado se sigue una implicancia potencial, no una conclusión de este estudio: si la
herramienta aumentara esa fracción, su efecto no se agotaría en quien la instala. Comprobarlo
exigiría medir el mercado y no al usuario, lo que excede el alcance de este trabajo.

La limitación principal se ha declarado a lo largo del artículo: la arquitectura no ha sido
implementada. El único resultado experimental disponible es de naturaleza diagnóstica y apunta
contra los propios datos: cinco algoritmos clasifican sin error la partición de prueba, lo que
revela que el conjunto es separable por construcción. Las métricas obtenidas sobre él no
estiman la calidad de la estimación de riesgo. Lo presentado es un diseño y un protocolo para
someterlo a prueba, no un sistema evaluado.

El trabajo futuro tiene un orden obligado, porque cada paso condiciona la interpretación del
siguiente. \emph{Primero}, medir qué características resultan legibles en fichas reales, con
especial atención al reparto entre plataformas grandes y vendedores independientes.
\emph{Segundo}, construir un conjunto de datos cuyas etiquetas provengan de una fuente
independiente del vector de entrada; sin él, cualquier métrica seguirá midiendo la regla que
generó esas etiquetas. \emph{Tercero}, implementar la arquitectura, incluido el reentrenamiento
del clasificador con la veracidad de dominio incorporada. \emph{Cuarto}, evaluar técnicamente
los componentes sobre ese conjunto independiente, con el análisis de ablación que establezca
cuánto aporta cada característica. \emph{Quinto}, y solo entonces, la evaluación con usuarios,
cuyas conclusiones dependen de que los cuatro pasos anteriores estén resueltos.

Una línea distinta es la detección del comportamiento fraudulento del vendedor, que exige
evidencia sobre operaciones consumadas y la colaboración de las plataformas; constituye una
recomendación para investigaciones posteriores antes que una extensión del sistema aquí
descrito.
